\cotefile{157-2}
\scanpage{275}
\mspage{16}

3. $\mathbb{\Delta}_2^+,\mathbb{\Delta}_2^-$, correspondent aux diagrammes
\diagram{\begin{tikzcd}[ampersand replacement=\&] a_0 \arrow[r] \& a_1 \& \unsure{\text{---}|\longrightarrow} \& a_2 \end{tikzcd}}
resp.\quad \deleted{\unsure{$a$}}
\diagram{\begin{tikzcd}[ampersand replacement=\&] a_0 \& a_1 \arrow[l] \& a_2 \arrow[l] \end{tikzcd}},
on peut faire une remarque similaire que pour $\mathbb{\Delta}_1^+,\mathbb{\Delta}_1^-$

4\unsure{,|.} $\Phi_1,\Psi_1$, ce sont respectivement les intervalles
\diagram{\begin{tikzcd}[ampersand replacement=\&,row sep=2pt,column sep=small]
{} \& {} \& {} \& {} \& {} \& {} \& 0 \\
\Phi_1\colon \& 0 \& 1 \arrow[l,"{\underline{u}_0}"'] \arrow[r,"{\underline{u}_1}"] \& 2 \& \text{ou} \& 1 \arrow[ur] \arrow[dr] \& {} \\
{} \& {} \& {} \& {} \& {} \& {} \& 2
\end{tikzcd}}
\diagram{\begin{tikzcd}[ampersand replacement=\&,row sep=2pt,column sep=small]
{} \& {} \& {} \& {} \& {} \& {} \& 0 \arrow[dl] \\
\Psi_1\colon \& 0 \arrow[r,"{\underline{u}_0}"] \& 1 \& 2 \arrow[l,"{\underline{u}_1}"'] \& \unsure{\text{ou}} \& 1 \& {} \\
{} \& {} \& {} \& {} \& {} \& {} \& 2 \arrow[ul]
\end{tikzcd}}

\deleted{Cette fois $\Phi$} Ils correspondent respectivement aux diagrammes
\diagram{\begin{tikzcd}[ampersand replacement=\&] a_0 \& a_1 \arrow[l,"{u_0}"'] \arrow[r,"{u_1}"] \& a_2 \end{tikzcd}}
et
\diagram{\begin{tikzcd}[ampersand replacement=\&] a_0 \arrow[r] \& a_1 \& a_2 \arrow[l] \end{tikzcd}}

\unsure{Les|les} catégories $\Phi_1,\Psi_1$ sont \deleted{\unsure{évidemment}} opposées l'une de l'autre, et ne sont pas isomorphes. Aussi n'y a-t-il \inserted{en général} aucun iso. entre $\underline{\mathrm{Ch}}_{\underline{\Phi}_1}(X)$ et $\underline{\mathrm{Ch}}_{\underline{\Psi}_1}(X)$,

Soit $I$ un intervalle orienté, et considérons l'\uline{intervalle opposé} $I^{\mathrm{opp}}$, ayant $=$ catégorie sous-jacente mais l'orientation opposée. \deleted{Alors} Comme la définition de la \deleted{\unsure{de}} catégorie des chemins de type $I$ ne fait pas intervenir l'orientation de $I$, on a donc un iso. canonique
\[ \underline{\mathrm{Ch}}_I(X) \unsure{\xrightarrow{\sim}|\simeq} \underline{\mathrm{Ch}}_{I^{\mathrm{opp}}}(X), \]
% archive sheet number: 274
